\scanpage{10}

structure de similitude \unsure{sss|ssi} $\exists\,a\in\Gamma\underline{\mathcal O}_{\unsure{\alpha|X}}^\ast$ (\uline{globalement}) telle que $Q'=aQ$, et cette $a$ est alors unique (\unsure{en quoi résume|en quoi réside|ce qui assure} l'unicité). \deleted{On appelle des} \inserted{Structure de similitude} fidèle : définie loc. par une $Q$ telle que $a\longmapsto aQ$ injectif (\deleted{cela ne dépend pas du choix de} du choix de \deleted{$Q$}, seulement de la structure de similitude).

1.2. \deleted{Soit}\ \deleted{\unsure{par|pour}} $P\subset\underline{\mathrm{Quad}}(E)$ est une structure de similitude, on peut regarder le sous-Module engendré, soit

\[\Delta=\Delta_P\subset\underline{\mathrm{Quad}}(E),\]

et on voit que

\[P\mathrel{\unsure{-|\longrightarrow|\longmapsto}}\Delta_P\]

donne une bijection entre l'ens. des structures de similitudes sur $E$, et l'ens des sous-Modules de $\underline{\mathrm{Quad}}(E)$ qui sont localement monogènes. Les \deleted{\unsure{formes q|formes}} structures \deleted{fidèl} de similitude \uline{fidèles} correspondent aux sous modules

\[\Delta\lhook\joinrel\longrightarrow\underline{\mathrm{Quad}}(E)\]

qui sont \uline{inversibles}.

\margin{\uline{NB} Dans le cas où on a un hom \uline{quelc.} $\Delta\longrightarrow\underline{\mathrm{Quad}}(E)$ (donc non \unsure{forc.} injectif) \unsure{s'in.}\ \deleted{\unsure{est}} \deleted{\unsure{fid}} définissant \deleted{une} \deleted{structure}}

1.3. Si $E$ est \inserted{\unsure{loc\textsuperscript{t}.|loc.}} \unsure{lib.|L.} de type fini, il en \deleted{\unsure{résulte}} est de $=$ de $\underline{\mathrm{Quad}}(E)\simeq\underline{\mathrm{Sym}}^2(E)$. Une structure de similitudes est dite \uline{strictement fidèle} (ou universellement fidèle) si l'hom. précédent est ``universellement injectif'', ou encore (ce qui revient au $=$) injectif sur chaque fibre résiduelle,
% archive sheet number: 9
